Automated code review and software quality evaluation are vital for modern software engineering, where systems evolve rapidly and continuous integration demands proactive quality assurance. However, conventional quality assurance practices rely primarily on modular static code analysis (SCA) operating at the file or class level. In distributed publish--subscribe architectures (ROS~2, DDS, MQTT), loose coupling among producers and consumers obscures indirect dependencies, creating an \textbf{Architecture-Code Gap}: a codebase may exhibit clean unit-level metrics while harboring severe architectural technical debt, such as single points of failure, co-located deployment bottlenecks, and transport contract mismatches. Furthermore, existing structural diagnostic frameworks operate open-loop---ranking components by scalar criticality without identifying the underlying architectural pathology, attributing it to an engineering role, or verifying remediation actions.

To address these challenges, we introduce \textbf{SaG-Prescribe}, an intelligent static system analysis framework that brings a deterministic \textbf{Diagnostic Pathway} to the forefront of automated code review and software quality evaluation. The framework unifies two complementary mechanisms: (1)~\textbf{Explainable Diagnostic Pathway \& Technical Debt Analysis:} It extracts a multi-layer heterogeneous dependency graph from Architecture-as-Code descriptors, synthesizes code-level SCA metrics into a \textbf{Code Quality Penalty (CQP)}, and computes a hierarchical \textbf{Reliability--Maintainability (RM)} quality attribution model grounded in ISO/IEC 25010 and ISO/IEC 25019 Quality-in-Use standards. It drives a catalog of nineteen formal publish--subscribe anti-patterns and bad smells with adaptive box-plot thresholds, providing explainable review decisions mapped to specific engineering roles (Reliability Engineers, SREs, Software Architects). (2)~\textbf{Prescriptive Refactoring with Counterfactual Verification:} It compiles diagnostic findings into three graph-mutation refactoring operators (logical topic splitting, physical anti-affinity reallocation, and transport QoS contract hardening). Crucially, every candidate refactoring is independently verified on a sandboxed counterfactual graph against a discrete-event cascade simulator, admitting only edits whose measured risk reduction exceeds simulator seed noise across a sweep of propagation thresholds ($\overline{\Delta I} > \kappa \sigma_{\text{seed}}$).

We evaluate SaG-Prescribe across eight benchmark scenarios spanning six architectural domains (autonomous vehicles, IoT smart cities, financial trading, healthcare, hub-and-spoke microservices, and hyper-scale enterprise systems, 29--3325 vertices), reporting empirical findings with full methodological transparency, including two calibration results a purely accuracy-framed evaluation would omit. On the diagnostic side, composite criticality, betweenness, and degree centrality all saturate around $\mathrm{NDCG@10} \approx 0.82$--$0.86$ and Cohen's $\kappa \approx 0.41$--$0.48$ against cascade-impact ground truth (mean Spearman $\rho = 0.485$ for the composite, app layer, versus $0.519$ for degree alone)---no single scalar ranking dominates, which is precisely why deterministic, role-attributed diagnosis (Reliability Engineer vs. SRE vs. Architect) is the pathway's value, not superior ranking accuracy. This re-measures, under an independent discrete-event cascade oracle, a criticality-ranking claim from our own prior work ($\rho = 0.94$ on a ten-application system against reachability loss); we report the full methodological delta (\S\ref{sec:intro:prior_work}, \S\ref{sec:results:attribution}). We also find that the anti-pattern catalog, while a highly sensitive component-level screen (mean recall $0.942$, precision $0.253$, implicating $94.8\%$ of components), agrees with ground-truth criticality at chance level (Cohen's $\kappa = -0.002$)---a consequence of $98.4\%$ of its 6329 findings (dominated by two edge-scoped patterns) not being attributable to any single component---and that the Availability dimension is measured on a degenerate substrate at this layer (zero predicted articulation points in 6 of 8 scenarios, SPOF detector $F_1 = 0.0$ throughout), motivating threshold recalibration, delta-aware gating, and layer-aware scoring. We also document a Simpson's paradox pooling effect across heterogeneous node types ($\rho = 0.028$ pooled vs. $\rho = 0.503$ on applications). On the prescriptive side, counterfactual verification admits 1128 of 1589 candidate refactorings ($71.0\%$), with anti-affinity reallocation demonstrating the highest survival rate ($77.7\%$), followed by transport QoS hardening ($58.0\%$) and topic splitting ($49.7\%$). Every scenario achieves a statistically significant net reduction in System Risk Index (Wilcoxon $W=0, p=0.0156, n=7$). The diagnostic review gate executes in 0.01--20.98~s, confirming its feasibility as a real-time, pre-commit automated code review check in CI/CD pipelines.
